% Notes of the Montreal seminar, 1970 (cote 129, typed part): archive cote 129, site page 18 (PDF page 19).
% typescript. Draft compared: none.
% AI-transcribed by Claude Opus 5.5 (claude-opus-5-5, adaptive thinking, effort high), Message Batches API, batch mtl129_full (msgbatch_01AGnDKHwy6ATCM7ehUa6263), 2026-11-13; not yet checked by eye.
\scanpage{129/18}
Par passage à la limite, les morphismes $V_n$ définissent un
morphisme additif
\[ V : \mathbb{W} \to \mathbb{W} \]
appelé \uline{Verschiebung} ou \uline{morphisme de décalage}.

2.2. \qquad Le \uline{morphisme de Frobenius} $F : \mathbb{W} \to \mathbb{W}$ est défini par passage
à la limite à partir des morphismes
\[ F_n : \mathbb{W}_n \to \mathbb{W}_n \qquad (x_1,x_2,\ldots,x_n) \to (x_1^p,x_2^p,\ldots,x_n^p) \ . \]

Ce morphisme n'est en général pas additif, mais par réduc-%
tion en caractéristique $p$ , il définit un \uline{homomorphisme d'anneau}
\[ F_{\mathbb{F}_p} : \mathbb{W}_{\mathbb{F}_p} \to \mathbb{W}_{\mathbb{F}_p} \qquad (\text{resp. } F_{n\,\mathbb{F}_p}) \]
où $\mathbb{F}_p$ désigne le corps premier $\mathbb{Z}/p\mathbb{Z}$ et $\mathbb{W}_{\mathbb{F}_p}$ désigne le schéma induit
sur $\mathbb{F}_p$ par changement de base. En effet, les deux lois de composition
sur $\mathbb{W}$ sont données par
\[ (x_1,x_2,\ldots,x_n,\ldots) + (y_1,\ldots,y_n,\ldots) = (S_1,S_2,\ldots,S_n,\ldots) \]
\[ (x_1,x_2,\ldots,x_n,\ldots) \cdot (y_1,\ldots,y_n,\ldots) = (P_1,\ldots,P_n,\ldots) \]
où les $S_i$ et $P_i$ sont des polynômes à coefficients entiers en
$(x_1,x_2,\ldots,x_i,\ y_1,y_2,\ldots,y_i)$ . Or, pour toute $\mathbb{F}_p$-algèbre, on a
\[ P(x^p,y^p) = [P(x,y)]^p \ . \]

\mspage{\hand{18}}
